\subsection{The sharp cost of one tanh row}
\label{sec:new-single-row-law}

The distinction between a value and a sample already appears in one
row. It is useful to state its complete query law before considering
depth. Auxiliary bit work is at least constant even when no input
coordinate is requested.

\begin{theorem}[One-row law]
\label{thm:new-single-row-law}
Let $x\in\{-1,1\}^n$, let $|b_0|\le1$, and let $w$ be a dyadic
row with $r=\sum_i|w_i|$. After an integer alias table has been built
from the row, an exact sign of mean
\[
 f(x)=\tanh\!\left(b_0+\sum_iw_ix_i\right)
\]
can be generated using at most $2r(1+r)$ expected source-coordinate
requests. Reading the entire input gives the alternative bound $n$.
Including random bits, row choices, and transcendental refinement,
the expected bit work is
\[
 O\!\left(B^4[1+\min\{n,r+r^2\}]\right),
\]
where $B$ includes the maximum encoding length of one coefficient,
the bias encoding, addresses, and $\log(2+r)$.

Conversely, when $n$ is a power of two and $r>0$ is dyadic, the
zero-bias row $w_i=r/n$ has worst-input expected distinct-probe
complexity $\Omega(\min\{n,r+r^2\})$, with an absolute implied
constant.
This lower bound permits arbitrary preprocessing independent of $x$.
The displayed row must, as usual, fit the stipulated coefficient grid.
For general widths, using the largest power of two at most $n$ changes
only the absolute constant.
\end{theorem}
\paragraph{Proof idea.}
Absolute-weight sampling turns the row into one unknown source mean, giving
the upper bound. For the lower bound, use a uniform row and compare its
sensitivity to the information obtained from distinct input probes.

\begin{proof}
If $r=0$, a known scalar coin suffices. Otherwise apply
Corollary~\ref{cor:comptanhrow} with $R=1$ and $h_i=x_i$: a source
request for coordinate $i$ is one input probe. This gives at most
$2r(1+r)$ expected input requests and $O((1+r)^2B^4)$ expected local
bit work, with coordinates selected through the integer alias table.
The dense alternative computes the dyadic preactivation after $n$
input probes and makes a known tanh coin by a certified uniform
comparison. Choose the cheaper alternative
using the known row norm. Both are exact and terminate almost surely.

For the lower bound write $M=n^{-1}\sum_ix_i$. If $0<r\le1$,
endpoint coupling (Section~\ref{sec:lower}) shows that some probe
occurs with probability at least the total variation distance
between the two endpoint output laws, namely $\tanh r\ge r/2$. Since
$r+r^2\le2r$, this proves the claim in this range.

Suppose $r\ge1$. Give the input signs independent common mean $t$,
and put $H(t)=\mathbb E_t\tanh(rM)$. At $t=0$, let
$Z=\sqrt n\,M$. Then
\[
 \mathbb E Z^2=1,\qquad \mathbb E Z^4\le3,\qquad
 \Pr(|Z|\ge1/2)\ge3/16.
\]
The last inequality is Paley--Zygmund applied to $Z^2$ at threshold
$1/4$. The finite-input score identity gives
\[
 H'(0)=\sqrt n\,\mathbb E\bigl[Z\tanh(rZ/\sqrt n)\bigr].
\]
For $u\ge0$, $\tanh u\ge\frac12\min\{u,1\}$: concavity gives
the assertion on $[0,1]$ using $\tanh1>1/2$, and monotonicity gives
it above one. Restricting the expectation to $|Z|\ge1/2$ therefore
gives
\[
 H'(0)\ge\frac{3}{128}\min\{r,\sqrt n\}.
\]
The transcript information inequality gives
\[
 \mathbb E_0Q\ge H'(0)^2
 \ge\frac9{16384}\min\{r^2,n\}
 \ge\frac1{65536}\min\{r+r^2,n\}.
\]
An input attaining at least this average expected cost exists. The
argument includes all preprocessing randomness in the independent
auxiliary tape and imposes no bound on preprocessing time.
\end{proof}

For fixed positive $r$, this law gives bounded expected input work
at arbitrarily large width. In contrast, the exact dyadic
preactivation of a row with every coefficient nonzero depends on
every input coordinate. Its deterministic evaluation must read all
of them: changing an unqueried sign changes the preactivation.
The network question is whether this saving survives composition.
